\documentclass[ngerman]{scrartcl}

%\usepackage{basis}

%\usepackage[draft]{basis}
%\usepackage[final]{basis}

%\usepackage[font=charter]{basis}
%\usepackage[font=hfold]{basis}
%\usepackage[font=mathpazo]{basis}
%\usepackage[font=original]{basis}
%\usepackage[font=times]{basis}
%\usepackage[font=wrong]{basis}

%\usepackage[hypercolor=green]{basis}

%\usepackage[hyperdriver=dvips]{basis}

%\usepackage[notitlepage]{basis}
%\usepackage[notitlepage, oneside]{basis}

%\usepackage[oneside]{basis}

%\usepackage[protokoll]{basis}

%\usepackage[pagestyle=beides]{basis}
%\usepackage[pagestyle=fuss]{basis}
\usepackage[pagestyle=fussseite]{basis}
%\usepackage[pagestyle=leer]{basis}
%\usepackage[notitlepage, pagestyle=leer]{basis}
%\usepackage[pagestyle=wrong]{basis}

%\usepackage[draft|final, font=charter|hfold|mathpazo|original|times, hypercolor=<color>, hyperdriver=<driver>, notitlepage, oneside, protokoll, pagestyle=beides|fuss|fussseite|leer]{basis}

%\title{Dokumenttitel}
\title[Kurzform]{Dokumenttitel}
\subtitle{Untertitel}

\version{Version 1.0}
\date{Datum}

\author{Autorname}

\strasse{Straße}
\plz{PLZ}
\ort{Ort}
\telefon{(030) 1\,22\,33\,44}
\handy{(0177) 1\,22\,33\,44}
\email{abc@abc.de}
\homepage{www.abc.de}
\adresszusatz{Adresszusatz}

\titelzusatz{Titelzusatz}
\logo{testlogo}

\usepackage{lipsum}

\begin{document}

	\maketitle

	\tableofcontents
	\cleardoublepage

	\section{Inhalte}
	\label{sec:sec}

	Text.

	Text \meta{Metatext} Text.

	\subsection{Bilder}
	\label{subsec:subsec}

	\begin{figure}[htb]
		\caption{Bild}
		\label{fig:fig}
		\includegraphics{testlogo}
	\end{figure}

	\subsection{Referenzen}

	Siehe fig \autoref{fig:fig}.

	Siehe sec \autoref{sec:sec}.

	Siehe subsec \autoref{subsec:subsec}.

	Siehe tab \autoref{tab:tab}.

	Siehe eqn \autoref{eqn:eqn}.

	\subsection{Symbole}

	\subsubsection{marvosym}

	Telefon \Telefon{}
	WashCotton \WashCotton{}
	Virgo \Virgo{}

	\subsubsection{wasysym}

	phone \phone{}
	virgo \virgo{}
	frownie \frownie{}

	\subsubsection{hologo}

	\hologo{XeLaTeX}
	\hologo{METAPOST}
	\hologo{LaTeXe}
	\hologo{pdfLaTeX}

	\subsection{Tabellen}

	\begin{longtabu}[l]{@{}lcr@{}}
		\toprule
		\captionabove{Tabelle über mehrere Seiten}\label{tab:tab}\\
		\midrule
		Zelle & Zelle & Zelleninhalt einer gewissen Länge\\
		Zelle & Zelleninhalt einer gewissen Länge & Zelle\\
		Zelle & Zelle & Zelle\\
		Zelle & Zelle & Zelleninhalt einer gewissen Länge\\
		Zelle & Zelleninhalt einer gewissen Länge & Zelle\\
		Zelle & Zelle & Zelle\\
		Zelle & Zelle & Zelleninhalt einer gewissen Länge\\
		Zelle & Zelleninhalt einer gewissen Länge & Zelle\\
		Zelle & Zelle & Zelle\\
		Zelle & Zelle & Zelleninhalt einer gewissen Länge\\
		Zelle & Zelleninhalt einer gewissen Länge & Zelle\\
		Zelle & Zelle & Zelle\\
		Zelle & Zelle & Zelleninhalt einer gewissen Länge\\
		Zelle & Zelleninhalt einer gewissen Länge & Zelle\\
		Zelle & Zelle & Zelle\\
		Zelle & Zelle & Zelleninhalt einer gewissen Länge\\
		Zelle & Zelleninhalt einer gewissen Länge & Zelle\\
		Zelle & Zelle & Zelle\\
		Zelle & Zelle & Zelleninhalt einer gewissen Länge\\
		Zelle & Zelleninhalt einer gewissen Länge & Zelle\\
		Zelle & Zelle & Zelle\\
		Zelle & Zelle & Zelleninhalt einer gewissen Länge\\
		Zelle & Zelleninhalt einer gewissen Länge & Zelle\\
		Zelle & Zelle & Zelle\\
		Zelle & Zelle & Zelleninhalt einer gewissen Länge\\
		Zelle & Zelleninhalt einer gewissen Länge & Zelle\\
		Zelle & Zelle & Zelle\\
		Zelle & Zelle & Zelleninhalt einer gewissen Länge\\
		Zelle & Zelleninhalt einer gewissen Länge & Zelle\\
		Zelle & Zelle & Zelle\\
		Zelle & Zelle & Zelleninhalt einer gewissen Länge\\
		Zelle & Zelleninhalt einer gewissen Länge & Zelle\\
		Zelle & Zelle & Zelle\\
		Zelle & Zelle & Zelleninhalt einer gewissen Länge\\
		Zelle & Zelleninhalt einer gewissen Länge & Zelle\\
		Zelle & Zelle & Zelle\\
		Zelle & Zelle & Zelleninhalt einer gewissen Länge\\
		Zelle & Zelleninhalt einer gewissen Länge & Zelle\\
		Zelle & Zelle & Zelle\\
		Zelle & Zelle & Zelleninhalt einer gewissen Länge\\
		Zelle & Zelleninhalt einer gewissen Länge & Zelle\\
		Zelle & Zelle & Zelle\\
		Zelle & Zelle & Zelleninhalt einer gewissen Länge\\
		Zelle & Zelleninhalt einer gewissen Länge & Zelle\\
		Zelle & Zelle & Zelle\\
		Zelle & Zelle & Zelleninhalt einer gewissen Länge\\
		Zelle & Zelleninhalt einer gewissen Länge & Zelle\\
		Zelle & Zelle & Zelle\\
		Zelle & Zelle & Zelleninhalt einer gewissen Länge\\
		Zelle & Zelleninhalt einer gewissen Länge & Zelle\\
		Zelle & Zelle & Zelle\\
		Zelle & Zelle & Zelleninhalt einer gewissen Länge\\
		Zelle & Zelleninhalt einer gewissen Länge & Zelle\\
		Zelle & Zelle & Zelle\\
		Zelle & Zelle & Zelleninhalt einer gewissen Länge\\
		Zelle & Zelleninhalt einer gewissen Länge & Zelle\\
		Zelle & Zelle & Zelle\\
		Zelle & Zelle & Zelleninhalt einer gewissen Länge\\
		Zelle & Zelleninhalt einer gewissen Länge & Zelle\\
		Zelle & Zelle & Zelle\\
		Zelle & Zelle & Zelleninhalt einer gewissen Länge\\
		Zelle & Zelleninhalt einer gewissen Länge & Zelle\\
		Zelle & Zelle & Zelle\\
		Zelle & Zelle & Zelleninhalt einer gewissen Länge\\
		Zelle & Zelleninhalt einer gewissen Länge & Zelle\\
		Zelle & Zelle & Zelle\\
		\bottomrule
	\end{longtabu}


	\subsection{Gleichungen}

	\begin{equation}
		a^2+b^2=c^2 \label{eqn:eqn}
	\end{equation}

	\begin{equation}
		a^2+b^2=c^2
	\end{equation}

	\begin{equation}
		a^2+b^2=c^2
	\end{equation}

	\begin{equation}
		a^2+b^2=c^2
	\end{equation}

	\begin{equation}
		a^2+b^2=c^2
	\end{equation}

	\begin{equation}
		a^2+b^2=c^2
	\end{equation}

	\begin{equation}
		a^2+b^2=c^2
	\end{equation}

	\begin{equation}
		a^2+b^2=c^2
	\end{equation}

	\begin{equation}
		a^2+b^2=c^2
	\end{equation}

	\begin{equation}
		a^2+b^2=c^2
	\end{equation}

	\begin{equation}
		a^2+b^2=c^2
	\end{equation}

	\begin{equation}
		a^2+b^2=c^2
	\end{equation}

	\begin{equation}
		a^2+b^2=c^2
	\end{equation}

	\begin{equation}
		a^2+b^2=c^2
	\end{equation}

	\begin{equation}
		a^2+b^2=c^2
	\end{equation}

	\begin{equation}
		a^2+b^2=c^2
	\end{equation}

	\begin{equation}
		a^2+b^2=c^2
	\end{equation}

	\subsection{Fußnoten}

	Text\footnote{Fußnote, die über mehrere Zeilen läuft, die über mehrere Zeilen läuft, die über mehrere Zeilen läuft, die über mehrere Zeilen läuft}.
	Text\footnote{Fußnote}.
	Text\footnote{Fußnote}.
	Text\footnote{Fußnote}.
	Text\footnote{Fußnote}.
	Text\footnote{Fußnote}\multiplefootnoteseparator\footnote{Fußnote}\footnote{Fußnote}\footnote{Fußnote}.
	Text\footnote{Fußnote}\footnote{Fußnote}.

	\lipsum[1-1]



	\addsec{Unnumeriert}

	\lipsum[1-1]

	\subsection*{Zwischenüberschrift}

	\lipsum[1-1]

	\subsubsection*{Unterzwischenüberschrift}

	\lipsum[1-1]



	\section{section}
	\lipsum[1-5]

	\section{section}
	\lipsum[1-5]

	\section{section}
	\lipsum[1-5]

	\section{section}
	\lipsum[1-5]

	\section{section}
	\lipsum[1-25]

\end{document}

